% ============================================================
%  Part XI — Nouns
% ============================================================
\gpart{XII}{Nouns}

% ------------------------------------------------------------
\gsec{Guessing the gender}

\gintro{Gender has to be learned with each noun, but a large number of endings predict it reliably. These lists cover most of everyday vocabulary.}

Gender is mostly arbitrary, but endings are a reliable shortcut. These rules
cover a large share of everyday vocabulary.

\tcap{Usually \textcolor{Masc}{\textbf{der}}}

\begin{flattbl}{colspec={X[0.8,l] X[1.1,l] X[0.8,l] X[1.1,l]}}
  -er    & der Lehrer   & -ling  & der Frühling \\
  -en    & der Wagen    & -ig    & der Honig \\
  -ismus & der Kapitalismus & -or & der Motor \\
  Days, months, seasons & der Montag & Weather & der Regen \\
\end{flattbl}

\tcap{Usually \textcolor{Fem}{\textbf{die}}}

\begin{flattbl}{colspec={X[0.8,l] X[1.1,l] X[0.8,l] X[1.1,l]}}
  -ung  & die Zeitung  & -heit & die Freiheit \\
  -keit & die Möglichkeit & -schaft & die Freundschaft \\
  -ion  & die Nation   & -tät  & die Universität \\
  -ie   & die Familie  & -ur   & die Kultur \\
  -e    & die Blume \eng{mostly} & -in & die Lehrerin \\
\end{flattbl}

\tcap{Usually \textcolor{Neut}{\textbf{das}}}

\begin{flattbl}{colspec={X[0.8,l] X[1.1,l] X[0.8,l] X[1.1,l]}}
  -chen & das Mädchen & -lein & das Fräulein \\
  -um   & das Zentrum & -ment & das Instrument \\
  -nis  & das Ergebnis & -tum & das Eigentum \\
  Infinitives as nouns & das Essen & Ge- & das Gebäude \\
\end{flattbl}

\begin{gnote}
\textit{-chen} and \textit{-lein} always win, whatever the base noun was:
\textit{die Frau} but \textit{\textbf{das} Fräulein}, \textit{der Hund} but
\textit{\textbf{das} Hündchen}. In a compound noun the gender comes from the
\textbf{last} element: \textit{die Tür} + \textit{der Griff} = \textit{der
Türgriff}.
\end{gnote}

% ------------------------------------------------------------
\gsec{Forming the plural}

\gintro{German has five plural patterns and no single rule for choosing between them, though gender and ending narrow the field considerably.}

There is no single plural rule — you learn it with the noun. These are the five
patterns and the tendencies that go with them.

\begin{gtbl}{colspec={X[0.75,l] X[1.15,l] X[1.1,l]}}
  Pattern & Typical of & Examples \\
  \en{-e} \eng{$\pm$ umlaut} & many masc. \& neut. & der Tag $\rightarrow$ die Tage \\
  \en{-er} \eng{$\pm$ umlaut} & short neuter nouns & das Kind $\rightarrow$ die Kinder \\
  \en{-(e)n} & almost all feminine & die Frau $\rightarrow$ die Frauen \\
  \en{-s}    & loanwords, abbreviations & das Auto $\rightarrow$ die Autos \\
  \en{--} \eng{no change} & -er, -en, -el endings & der Lehrer $\rightarrow$ die Lehrer \\
  \en{-nen}  & feminine -in & die Lehrerin $\rightarrow$ die Lehrerinnen \\
\end{gtbl}

\begin{gnote}
The plural article is always \textcolor{Plur}{\textbf{die}}, whatever the
singular gender was. And in the dative plural the noun picks up an extra
\textit{-n} unless it already ends in \textit{-n} or \textit{-s}.
\end{gnote}

\tcap{Umlaut in the plural}

\begin{flattbl}{colspec={X[1,l] X[1,l] X[1,l]}}
  der Mann $\rightarrow$ die Männer & die Hand $\rightarrow$ die Hände & das Haus $\rightarrow$ die Häuser \\
  der Stuhl $\rightarrow$ die Stühle & die Stadt $\rightarrow$ die Städte & das Buch $\rightarrow$ die Bücher \\
  der Vater $\rightarrow$ die Väter & die Mutter $\rightarrow$ die Mütter & das Wort $\rightarrow$ die Wörter \\
\end{flattbl}

% ------------------------------------------------------------
\gsec{Weak nouns (n-declension)}

\gintro{A closed group of masculine nouns adds \textit{-n} or \textit{-en} in every case except the nominative singular. They are easy to spot once you know the group.}

A closed group of masculine nouns that take \textbf{-n} or \textbf{-en} in every
case except the nominative singular.

\begin{gendertbl}{}
  Case & der Junge & der Herr & der Name & der Student \\
  Nominativ & der Junge  & der Herr   & der Name   & der Student \\
  Akkusativ & den Jungen & den Herrn  & den Namen  & den Studenten \\
  Dativ     & dem Jungen & dem Herrn  & dem Namen  & dem Studenten \\
  Genitiv   & des Jungen & des Herrn  & des Namens & des Studenten \\
\end{gendertbl}

\tcap{Who belongs to the group}

\begin{flattbl}{colspec={X[1,l] X[2,l]}}
  Masc. in -e   & der Junge, der Kunde, der Kollege, der Neffe, der Löwe \\
  People in -ent, -ant, -ist, -at & der Student, der Polizist, der Soldat \\
  Nationalities & der Franzose, der Grieche, der Türke, der Ire \\
  Odd ones out  & der Herr \eng{-n sg., -en pl.}, der Mensch, der Nachbar, der Bauer \\
  Mixed         & der Name, der Gedanke, der Glaube \eng{add -s in the genitive} \\
\end{flattbl}

\begin{gnote}
The classic error is \textit{Ich sehe den Student}. It has to be \textit{den
Student\textbf{en}}. If a masculine noun ends in \textit{-e} and refers to a
person or an animal, assume it is weak.
\end{gnote}

% ------------------------------------------------------------
\gsec{Compound nouns}

\gintro{German joins nouns together rather than stringing them out with prepositions. The result is one word, however long.}

German builds long nouns by stacking short ones. Read them right to left.

\begin{gtbl}{colspec={X[1.35,l] X[0.85,l] X[0.8,l]}}
  Compound & Parts & Gender from \\
  die Haustür        & Haus + Tür        & die Tür \\
  der Fußballspieler & Fußball + Spieler & der Spieler \\
  das Krankenhaus    & krank + Haus      & das Haus \\
  die Bahnhofstraße  & Bahnhof + Straße  & die Straße \\
  die Geschwindigkeitsbegrenzung & Geschwindigkeit + Begrenzung & die Begrenzung \\
\end{gtbl}

\begin{gnote}
The last element carries the gender, the plural and the case ending. Everything
in front of it is just description. A linking \textit{-s} or \textit{-n} often
appears at the seam (\textit{Arbeit\textbf{s}platz},
\textit{Kranke\textbf{n}haus}); it follows historical patterns rather than a
productive rule.
\end{gnote}

\tcap{Capitalisation}

\begin{flattbl}{colspec={X[1,l] X[1.9,l]}}
  Every noun & \textit{der Hund}, \textit{die Freiheit}, \textit{das Essen} \\
  Nouns made from other words & \textit{das Laufen}, \textit{das Gute}, \textit{etwas Neues} \\
  Formal \textit{Sie}, \textit{Ihr} & always capitalised \\
  Informal \textit{du}, \textit{ihr} & lower case by default; capitalised \textit{Du}/\textit{Ihr} allowed again in letters \eng{since 2006} \\
\end{flattbl}

% ------------------------------------------------------------
\gsec{Countries and nationalities}

\gintro{Most country names take no article, and those are the easy ones. A
minority carry an article, and that choice changes the preposition you use with
them.}

\tcap{With an article}

\begin{gtbl}{colspec={X[1.05,l] X[0.85,l] X[1.1,l]}}
  Country & Gender & In / to \\
  die Schweiz     & feminine & in der / in die Schweiz \\
  die Türkei      & feminine & in der / in die Türkei \\
  die Niederlande & plural   & in den / in die Niederlande \\
  die USA         & plural   & in den / in die USA \\
  der Iran        & masculine & im / in den Iran \\
  das Vereinigte Königreich & neuter & im / in das \dots \\
\end{gtbl}

\begin{gnote}
\textit{nach} is used for places without an article, and also in
\textit{nach Hause}. Places carrying an article take \textit{in} plus the
accusative for movement and the dative for position. People and businesses take
\textit{zu}: \textit{zum Arzt}. Coming from somewhere is \textit{aus}:
\textit{aus Deutschland}, \textit{aus der Schweiz}.
\end{gnote}

\tcap{Nationalities and languages}

\begin{gtbl}{colspec={X[0.9,l] X[1.05,l] X[1.05,l]}}
  Country & Person & Adjective \\
  Deutschland & der Deutsche / die Deutsche & deutsch \\
  Frankreich  & der Franzose / die Französin & französisch \\
  England     & der Engländer / die Engländerin & englisch \\
  Spanien     & der Spanier / die Spanierin & spanisch \\
  die Türkei  & der Türke / die Türkin & türkisch \\
\end{gtbl}

\begin{gnote}
Nationality nouns and language names are capitalised; only the adjective is
lower case. \textit{Ich bin \textbf{Deutscher}}, \textit{Ich spreche
\textbf{Deutsch}}, \textit{auf \textbf{Deutsch}} — but \textit{ein
\textbf{deutsches} Auto}.
\end{gnote}
